\documentclass[../main.tex]{subfiles}
\begin{document}
\chapter{The Second Variation}
\section{Defining the Second Variation}
The \textit{second variation} of a functional is analogous to the second derivative of a function.
\begin{remark}[Recap]
  Recall that for a function $f$:
  \[
    f(\vec{x} + \varepsilon \vec{\xi}) - f(\vec{x}) = \varepsilon \vec{\xi} \cdot \nabla f(\vec{x}) + \frac{1}{2} \varepsilon^2 \xi_i \xi_j H_{i j}(\vec{x}) + O(\varepsilon^3)
  \]
  It follows that the point $\vec{x} = \vec{a}$ is a local minimum if $H_{i j}(\vec{a})$ is \textit{positive definite}, i.e. $\xi_i\xi_j H_{i j}(\vec{a}) > 0\ \forall \vec{\xi} \neq \vec{0}$, or equivalently, all eigenvalues of $H_{i j}(\vec{a})$ are positive.
\end{remark}
Consider the general functional:
\[
  F[y] = \int_{\alpha}^{\beta} f(y, y', x) \d{x}
\]
where $y \in \mathscr{C}$ and:
\[
  \mathscr{C} = \{(y: [\alpha, \beta] \to \R) : y(\alpha) = a, y(\beta) = b, \text{$y$ is twice continuously differentiable}\}
\]
Consider a perturbation $y + \varepsilon \xi \in \mathscr{C}$ with $\xi(\alpha) = \xi(\beta) = 0$ so that the boundary conditions still hold.
Expanding to second order:
\begin{align*}
  f(y + \varepsilon \xi, y' + \varepsilon \xi', x) &= f(y, y', x) + \varepsilon \xi\left[\pderiv{f}{y} - \deriv{}{x} \pderiv{f}{y'}\right] + \varepsilon \deriv{}{x}\left(\xi \pderiv{f}{y'}\right) \\
                                                   &\quad + \frac{1}{2} \varepsilon^2 \left[\xi^2 \pderiv[2]{f}{y} + 2 \xi \xi' \frac{\partial^2 f}{\partial y \partial y'} + {\varepsilon'}^2 \pderiv[2]{f}{{\xi'}}\right]
\end{align*}
The first two terms follow from how we derived the Euler-Lagrange equation, just not inside of an integral this time.
Upon substituting this back in:
\begin{align*}
  F[y + \varepsilon \xi] - F[y] &= \varepsilon \int_{\alpha}^{\beta} \xi(x) \frac{\delta F[y]}{y(x)} \d{x} + \varepsilon^2 \delta^2 F[y, \xi] + O(\varepsilon^3)
\end{align*}
where the boundary terms vanish because $\xi(\alpha) = \xi(\beta) = 0$.
We define:
\[
  \delta^2 F[y, \xi] \equiv \frac{1}{2} \int_{\alpha}^{\beta} \left[\xi^2 \pderiv[2]{f}{y} + 2\xi \xi' \frac{\partial^2 f}{\partial y \partial y'} + {\xi'}^2 \pderiv{f}{{\xi'}}\right] \d{x}
\]
Now,  utilising that $2\xi \xi' = \deriv{}{x}(\xi^2)$, we can integrate by parts:
\[
  \delta^2 F[y, \xi] = \int_{\alpha}^{\beta} \left[\xi^2 \left(\pderiv[2]{f}{y} - \deriv{}{x} \frac{\partial^2 f}{\partial y \partial y'}\right) + {\xi'}^2 \pderiv[2]{f}{{y'}}\right] \d{x}
\]
again utilising the boundary conditions.
\begin{remark}[Recap]
  Recall that if $H(\vec{x})$ is positive definite for all $\vec{x}$, then $f(\vec{x})$ is convex and so any stationary point must be a global minimum.
\end{remark}
We can formulate an analogous statement to the above but instead for functionals.
\begin{theorem}
  If $\delta^2 F[y, \xi] \geq 0$ for all $y \in \mathscr{C}$ and allowable $\xi(x)$, then $y = y_0$ is a global minimum of $F$ within $\mathscr{C}$ if $y_0 \in \mathscr{C}$ satisfies the Euler-Lagrange equation.
\end{theorem}
\begin{remark}[Note]
  \begin{itemize}
    \item When we refer to ``allowed'' $\xi(x)$, we mean that it obeys the boundary conditions set up previously.
    \item For a function to be convex, we require that its domain is also convex.
      This is true for $\mathscr{C}$ as if $y_1, y_2 \in \mathscr{C}$ then so is $ty_1 + (1 - t)y_2$ for all $t \in [0, 1]$.
  \end{itemize}
\end{remark}
\begin{example}[Geodesics Revisted]
  Recall that we used the functional $L[y]$ for the length of the curve:
  \[
    L[y] = \int_{\alpha}^{\beta} \sqrt{1 + (y')^2} \d{x}
  \]
  There is no explicit $y$ dependence within the integrand $f(y, y', x) = \sqrt{1 + (y')^2}$ and so:
  \[
    \pderiv[2]{f}{y} = \frac{\partial^2 f}{\partial y \partial y'} = 0
  \]
  and we can find:
  \[
    \pderiv[2]{f}{{y'}} = (1 + (y')^2)^{-\frac{3}{2}}
  \]
  Therefore the second variation is given by:
  \[
    \delta^2 F[y, \xi]a = \frac{1}{2} \int_{\alpha}^{\beta} {\xi'}^2 (1 + (y')^2)^{-\frac{3}{2}} \d{x} \geq 0\ \forall y, \xi
  \]
  Thus any solution of the Euler-Lagrange equation (which we saw was a straight line) is a global minimum of the distance between two points.
\end{example}
\section{Classifying Optimal Solutions}
If $y_0 \in \mathscr{C}$ satisfies the Euler-Lagrange equation then:
\begin{itemize}
  \item If $\delta^2 F[y_0, \xi] > 0$ for all allowed $\xi(x)$ that are not identically 0, then $y_0$ is a local minimum.
    Conversely, we can conclude that $y_0$ is not a local minimum if we can find any $\xi(x)$ such that $\delta^2 F[y_0, \xi] < 0$.
  \item If $\delta^2 F[y_0, \xi] < 0$ for all allowed $\xi(x)$ that are not identically 0, then $y_0$ is a local maximum.
    Conversely, we can conclude that $y_0$ is not a local maximum if we can find any $\xi(x)$ such that $\delta^2 F[y_0, \xi] > 0$.
\end{itemize}
\begin{example}[Geodesics on a Circle]
  On Example Sheet 2, you will show that a geodesic between two points $A, B$ on a circle is a segment of the unique great circle that passes through $A, B$.

  For this unique great circle, there is two segments we could choose, and so we calculate the second variation to find that the shorter one has $\delta^2 F > 0\ \forall \xi \neq 0$ and so it must be a local minimum, and in turn is a global minimum.

  The longer of the two paths can be shown to be non-minimal as we can find $\xi$ such that $\delta^2 F < 0$.

  In the case where $A$ and $B$ are antipodal points, then $\delta^2 F \geq 0\ \forall \xi$ which is not covered by our criteria.
\end{example}
\subsection{Legendre Condition}
\begin{proposition}[Lengendre Condition]
  Given any $y_0 \in \mathscr{C}$ satisfying the Euler-Lagrange equations, a necessary condition for $y_0$ to be a local minimum is:
\[
    \at{\pderiv[2]{f}{{y'}}}{y = y_0} \geq 0\ \forall x \in [\alpha, \beta]
  \]
\end{proposition}
\begin{proof}
  Suppose that $\at{\pderiv[2]{f}{{y'}}}{y = y_0} < 0$ at some $x_0 \in [a, b]$, then we can choose $\xi$ such that $\xi$ is non-zero only near $x_0$ and $\xi'$ is large with $\xi$ small to make $\delta^2 F[y_0, \xi] < 0$ and so $y_0$ cannot be a local minimum.

  Thus taking the contrapositive of the above, we obtain the result.
\end{proof}
\begin{example}
  Consider the functional:
  \[
    F[y] = \int_{-1}^{1} x\sqrt{1 + (y')^2} \d{x}
  \]
  We then have:
  \[
    \pderiv[2]{f}{{y'}} = x(1 + (y')^2)^{-\frac{3}{2}}
  \]
  which violates the Legendre condition for $x < 0$.
  Thus any solution of the Euler-Lagrange equations is \textbf{not} a local minimum.
\end{example}
Let $y_0 \in \mathscr{C}$ satisfy the Euler-Lagrange equation and write:
\begin{equation}
  \delta^2 F[y_0, \xi] = \frac{1}{2} \int_{\alpha}^{\beta} [\rho(x) {\xi'}^2 + \sigma(x)\xi^2] \d{x} \label{2ndVar}
\end{equation}
where we define:
\[
  \rho = \pderiv[2]{f}{{y'}} \text{ and } \sigma = \left[\pderiv[2]{f}{y} - \deriv{}{x} \frac{\partial^2 f}{\partial y \partial y'}\right]
\]
We can then reformulate the Legendre condition as $\rho \geq 0$.

\begin{proposition}
  A sufficient conditions that $y_0$, a solution of the Euler-Lagrange equation, to be minimum is $\rho > 0$ and $\sigma \geq 0$.
\end{proposition}
\begin{proof}
  Clearly $\delta^2 F \geq 0\ \forall \xi$.
  Now suppose that $\delta^2 F = 0$ for some $\xi$.
  We then must have $\xi' = 0$ on $[\alpha , \beta]$, and thus $\xi$ is constant on $\alpha, \beta$.
  Since $\xi \in \mathscr{C}$, $\xi(\alpha) = \xi(\beta) = 0$ and so $\xi$ is identically 0 which is a contradiction and so $\delta^2 F > 0\ \forall \xi$.
\end{proof}
\begin{example}[Brachistochrone Revisited]
  Recall that we were aiming to minimise the functional:
  \[
    T[y] = \int_{0}^{x_B} f \d{x} \text{ where } f= \sqrt{\frac{1 + {y'}^2}{- y}}
  \]
  We need $y$ to be negative and so we impose this additional restriction on $\mathscr{C}$.

  Now assume that $y$ satisfies the Euler-Lagrange equations (which we saw is a cycloid).
  We see that:
  \[
    \pderiv{f}{{y'}} = \frac{y'}{\sqrt{-y(1 + {y'}^2)}},\ \pderiv{f}{y} = - \frac{f}{2y}
  \]
  and:
  \[
    \rho = \pderiv[2]{f}{{y'}} = \frac{1}{\sqrt{(-y)(1 + {y'}^2)^3}} > 0
  \]
  so the first condition is satisfied.

  To find $\sigma$ we need to calculate:
  \[
    \frac{\partial^2 f}{\partial y \partial y'} = - \frac{1}{y} \pderiv{f}{y'} \text{ and } \pderiv[2]{f}{y} = \frac{3f}{4y^2}
  \]
  and so $\sigma$ is:
  \begin{align*}
    \sigma &= \pderiv[2]{f}{y} - \deriv{}{x}\left(-\frac{1}{2y} \pderiv{f}{{y'}}\right) \\
           &= \pderiv[2]{f}{y} - \frac{y'}{2y^2} \pderiv{f}{{y'}} + \frac{1}{2y} \pderiv{f}{y} \text{ by using the EL equations} \\
           &= \frac{3f}{4y^2} - \frac{{y'}^2}{2y^2\sqrt{(-y)(1 + {y'}^2)}} - \frac{f}{4y^2} \\
           &= \frac{1}{2\sqrt{(-y)(1 + {y'}^2)}} > 0
  \end{align*}
  Thus $\rho > 0$ and $\sigma > 0$ so $y$ (a cycloid) is a local minimum of $T$.
\end{example}
We can also integrate \cref{2ndVar} by parts to yield:
\begin{align*}
  \delta^2 F[y_0, \xi] &= \frac{1}{2} \int_{\alpha}^{\beta} \xi(x) \mathcal{L}\xi(x) + \frac{1}{2}\underbrace{\eval{\rho \xi \xi'}{\alpha}{\beta}}_{= 0 } \d{x} \\
                       &= \frac{1}{2}\int_{\alpha}^{\beta} \xi \mathcal{L}\xi(x) \d{x}
\end{align*}
where $\mathcal{L}$ is the differential operator defined as:
\[
  \mathcal{L}\xi \equiv -\deriv{}{x}\left(\rho \deriv{\xi}{x}\right) + \rho \xi
\]
\begin{proposition}
  Consider the Sturm Louisville problem:
  \[
    \mathcal{L} \eta = \lambda w(x) \eta \text{ and } \eta(\alpha) = \eta(\beta) = 0
  \]
  where $w(x) > 0$ on $(\alpha, \beta)$.
  Then $y_0$ is a local minimum if and only if all eigenvalues are positive.
  s roblem admits a negative eigenvalue $\lambda$ then $y_0$ \textbf{is not} a local minimum of $F$.
\end{proposition}
\begin{proof}
  \begin{proofdirection}{Assume $y_0$ is a local minimum}
    Suppose for contradiction that there is an eigenfunction $\eta$ of $\mathcal{L}$ with negative eigenvalue $\lambda$.
    We then have:
    \begin{align*}
      \delta^2 F[y_0, \eta] &= \frac{1}{2} \int_{\alpha}^{\beta} \eta \mathcal{L} \eta \d{x} \\ \
                            &= \frac{\lambda}{2} \int_{\alpha}^{\beta} w \eta^2 \d{x} < 0
    \end{align*}
    and so $y_0$ cannot be a local minimum.
    Thus all eigenvalues must be positive.
  \end{proofdirection}
  \begin{proofdirection}{Assume all eigenvalues are positive}
    \nonexaminable
    Let $\lambda_n$ be the $n$-th eigenvalue of $\mathcal{L}$ and $y_n$ the $n$-th eigenfunction of $\mathcal{L}$.
    From IB Methods, the eigenfunctions $\{y_n\}$ form a complete space for the set of allowed $\eta$, i.e. we can write:
    \[
      \eta(x) = \sum_{n} a_n y_n(x)
    \]
    for some coefficients $\{a_n\}$.
    We can also choose $\{y_n\}$ to be orthogonal, i.e.:
    \[
      \int_{\alpha}^{\beta} w y_m y_n \d{x} = 0 \text{ for } m \neq n
    \]
    and so:
    \begin{align*}
      \delta^2 F[y_0, \eta] &= \frac{1}{2} \sum_{m, n} a_m a_n \lambda_n \int_{\alpha}^{\beta} w y_m y_n \d{x} \\
                            &= \frac{1}{2} \sum_{m} \lambda_m a^2_m \int_{\alpha}^{\beta} w y^2_m \d{x} \geq 0
    \end{align*}
    which vanishes if and only if $a_m = 0\ \forall m$.
    Thus $\delta^2 F > 0$ for $\eta \centernot\equiv 0$.
  \end{proofdirection}
\end{proof}
\begin{example}
  Consider the functional: \label{2ndVarExample}
  \[
    F[y] = \frac{1}{2} \int_{0}^{a} ({y'}^2 - y^2) \d{x}
  \]
  The Euler-Lagrange equations tell us:
  \[
    - y - \deriv{}{x} y' = 0 \iff y'' + y = 0
  \]
  which we can easily solve to yield a solution $y_0$.

  Expanding around this solution $y_0$ we have:
  \[
    \delta^2 F = \frac{1}{2} \int_{0}^{a} ({\xi'}^2 - \xi^2) \d{x}
  \]
  which tells us that $\rho = 1$ and $\sigma = -1$.
  Thus the associated Sturm Louisville problem is:
  \[
    -\deriv[2]{\eta}{x} - \eta = \lambda w \eta \text{ and } \eta(0) = \eta(a) = 0
  \]
  If we choose $w(x) \equiv 1$ then:
  \[
    \eta'' + (\lambda + 1)\eta = 0 \implies \eta = A \sin(x\sqrt{\lambda + 1}) + B \cos(x\sqrt{\lambda + 1})
  \]
  For a non-trivial $\eta(x)$, we need $\sqrt{\lambda + 1} a = n\pi$ for some $n \in \Z$ and so we must have:
  \[
    \eta(x) = A \sin \left(\frac{n \pi x}{a}\right)
  \]
  for $n \in \N$ as $n = 0$ is trivial and $n < 0$ can be absorbed by changing the sign of $A$.

  The eigenvalues are:
  \[
    \lambda_n = \frac{n^2 \pi^2}{a} - 1
  \]
  We then see that:
  \begin{align*}
    \text{local min} \iff \lambda_n > 0\ \forall n \iff a < \pi \\
    \text{not a local min} \text{ if } \exists n \text{ s.t. } \lambda_n < 0 \text{ if } a > \pi
  \end{align*}
\end{example}
\subsection{Jacobi Condition}
\nonexaminable
Ley $y_0 \in \mathscr{C}$ satisfy the Euler Lagrange equation.
Since $\xi(\alpha) = \xi(\beta) = 0$ we have:
\[
  0 = \int_{\alpha}^{\beta} (\varphi \xi^2)' \d{x} = \int_{\alpha}^{\beta} (2 \varphi \xi \xi' + \varphi' \xi^2) \d{x}
\]
for any differentiable $\varphi$.
Adding this to $\delta^2 F[y_0, \xi]$ we have:
\[
  \delta^2 F[y_0, \xi] = \frac{1}{2} \int_{\alpha}^{\beta} (\rho {\xi'}^2 + 2 \varphi \xi \xi' + (\sigma + \varphi')\xi^2) \d{x}
\]
Assuming that $\rho > 0$ for $\alpha \leq x \leq \beta$, we can complete the square inside the integral as:
\[
  \delta^2 F[y_0, \xi] = \frac{1}{2} \int_{\alpha}^{\beta} \rho\left(\xi' + \frac{\varphi}{\rho}\xi\right)^2 + \left(\sigma + \varphi' - \frac{\varphi^2}{\rho}\right)\xi^2 \d{x}
\]
We want to choose a $\varphi$ to satisfy the \textit{Ricatti equation}:
\begin{equation}
  \varphi' + \sigma = \frac{\varphi^2}{\rho} \label{ricatti}
\end{equation}
then $\delta^2 F[y_0, \xi] \geq 0$ and vanishes if and only if $\xi' + \frac{\varphi}{\rho}\xi = 0$ i.e. if and only if $\xi$ is of the form:
\[
  \xi(x) = C \exp\left[-\int_{\alpha}^{x} \frac{\varphi(s)}{\rho(s)} \d{s}\right]
\]
but $\xi(\alpha) = 0 \implies C = 0$ and so $\xi(x) \equiv 0$.
Thus $\delta^2 F[y_0, \xi] > 0$ for all non-trivial allowed $\xi$, and so $y_0$ is a local minimum.
So in summary, $y_0$ is a local minimum if there exists a solution to \cref{ricatti} on $[\alpha, \beta]$.

We see that \cref{ricatti} is a nonlinear ODE, however we can convert to a linear second order ODE by performing the substitution $\phi = - \rho \frac{u'}{u}$ for $u \neq 0$.
This transforms \cref{ricatti} into:
\begin{align*}
  -\left(\rho \frac{u'}{u}\right)' + \sigma &= \rho \left(\frac{u'}{u}\right)^2 \\
   \iff - \frac{(\rho u')'}{u} + \cancel{\rho \frac{{u'}^2}{u^2}} + \sigma &= \cancel{\rho \frac{{u'}^2}{u^2}} \\
   \iff -(\rho u')' + \sigma u &= 0
\end{align*}
which is the \textit{Jacobi accessory equation}.
If there exists a solution to this with $u \neq 0$ on $[\alpha, \beta]$ then we obtain a solution $\varphi$ of \cref{ricatti} and so $y_0$ is a local minimum.
\begin{remark}
  Since the Jacobi accessory equation is a linear ODE, it always has a solution but not necessarily one with $u \neq 0$.
\end{remark}
\begin{example}
  Consider again the functional from \cref{2ndVarExample}.
  Recall that we had $\rho = 1$ and $\sigma = -1$, thus the Jacobi accessory equation is:
  \[
    -u'' - u = 0 \implies u = A \sin (x - x_0)
  \]
  Thus the solutions oscillate with period $2\pi$  and so we can find a solution with $u = 0$ on $[0, a]$.
  Therefore the solution of the Euler-Lagrange equation is a local minimum if $a < \pi$ as we found with the Sturm-Louisville approach in \cref{2ndVarExample}.
  Unlike the Sturm-Louisville approach, this method does not let us prove that the solution is not a local minimum for $a > \pi$.
\end{example}
\end{document}
